\documentclass[11pt]{article}\usepackage{layout}\geometry{paperwidth=420mm,paperheight=297mm,margin=0mm}\begin{document}
\begin{canvas}
\pdfbookmark[1]{图 2-1}{fig1}\Box{17}{10}{386}{\includegraphics[width=386mm]{.build/assets/01-global-current.pdf}}
\Box{17}{237}{185}{\fontsize{12.5}{18}\selectfont\color{teal}\sffamily\bfseries 先看参与在哪里，再看平台看见了什么}
\Box{17}{248}{185}{\fontsize{11.1}{18.5}\selectfont\color{ink}美国、印度、中国和巴西位于本图活跃规模前列，欧洲则呈现多个相邻参与节点。地图提供合作伙伴的空间线索，前 20 排行补充可比的数量读数；二者共同呈现参与基础，不能据此判断技术质量、AI 研发实力或实际社会成效。}
\Box{218}{237}{185}{\fontsize{12.5}{18}\selectfont\color{teal}\sffamily\bfseries 中国观察需要容纳本土平台}
\Box{218}{248}{185}{\fontsize{11.1}{18.5}\selectfont\color{ink}Gitee、AtomGit 等平台使一部分中国参与力量进入视野。全球排行仍统一采用 GitHub 口径；本土平台补充另行展示，不与全球排行混排。五域趋势中的中国约 350 万采用独立汇总序列，既不替代地图的 250.48 万，也不由旧补充估计直接相加得到。}
\Rule{17}{285}{403}{line}\Small{17}{289}{350}{GOSD 2026 / 全球开源发展报告 · 精要版}\Box{385}{288}{18}{\raggedleft\digits\fontsize{11}{13}\selectfont 01}
\end{canvas}
\begin{canvas}
\pdfbookmark[1]{图 2-2}{fig2}\Box{17}{10}{386}{\includegraphics[width=386mm]{.build/assets/02-trends-2030.pdf}}
\Box{17}{237}{185}{\fontsize{12.5}{18}\selectfont\color{teal}\sffamily\bfseries 增长潜力与既有规模，是两种不同的优势}
\Box{17}{248}{185}{\fontsize{11.1}{18.5}\selectfont\color{ink}按修订后的历史序列，印度和南方六国代表组的年均增速较高，中国也持续扩大参与基础。若历史增速延续，到 2030 年两组成长群体相对于美国、欧洲的规模比例可能提高。但欧洲、美国仍具有较大的绝对规模；不能把相对增速更快写成整体能力全面领先。}
\Box{218}{237}{185}{\fontsize{12.5}{18}\selectfont\color{teal}\sffamily\bfseries 位次会随假设改变，预测应成为持续观察的问题}
\Box{218}{248}{185}{\fontsize{11.1}{18.5}\selectfont\color{ink}在历史增速延续情景下，印度 2030 年的活跃规模高于中国；在逐步放缓情景下，两者接近且次序不同。更稳妥的判断是合作版图可能扩大，具体位次尚有不确定性。累计用户与当期活跃的增长也不完全同步，账号规模扩大仍须转化为持续贡献。}
\Rule{17}{285}{403}{line}\Small{17}{289}{350}{GOSD 2026 / 全球开源发展报告 · 精要版}\Box{385}{288}{18}{\raggedleft\digits\fontsize{11}{13}\selectfont 02}
\end{canvas}
\begin{canvas}
\pdfbookmark[1]{图 2-3}{fig3}\Box{17}{10}{386}{\includegraphics[width=386mm]{.build/assets/03-china-regions.pdf}}
\Box{17}{237}{185}{\fontsize{12.5}{18}\selectfont\color{teal}\sffamily\bfseries 头部集聚与多地参与同时存在}
\Box{17}{248}{185}{\fontsize{11.1}{18.5}\selectfont\color{ink}京沪粤合计约 131.73 万活跃开发者，占本图全国规模的 52.6\%；第 4—10 位地区合计约占 36.0\%。头部之外的参与并非可以忽略的余量。累计规模、当期活跃与净增量并列，帮助地方决策者区分已有基础、当期参与和变化速度。}
\Box{218}{237}{185}{\fontsize{12.5}{18}\selectfont\color{teal}\sffamily\bfseries 增量结构为区域合作提供进一步线索}
\Box{218}{248}{185}{\fontsize{11.1}{18.5}\selectfont\color{ink}台湾与浙江合计约占当期活跃规模的 18.9\%，却贡献约 23.3\% 的净增量；这一差异值得在后续项目与人才分析中追踪。净增量是两期人数之差，不是首次参与人数，也不表示人员迁入。省级观察不能直接替代杭州等城市的具体证据。}
\Rule{17}{285}{403}{line}\Small{17}{289}{350}{GOSD 2026 / 全球开源发展报告 · 精要版}\Box{385}{288}{18}{\raggedleft\digits\fontsize{11}{13}\selectfont 03}
\end{canvas}
\end{document}